\documentclass[twocolumn]{article}
\usepackage{xeCJK} % 支持中文
\usepackage[left=2cm,right=2cm,top=2cm,bottom=2cm]{geometry} % 设置页边距
\usepackage{abstract} % 定制摘要

\usepackage{indentfirst} % 设置首行缩进

% 原模板使用 SimSun/SimHei；macOS 无此二字体，改用系统自带宋体/黑体
\setCJKmainfont{Songti SC}[BoldFont = Heiti SC]
\setCJKsansfont{Heiti SC}
\setmainfont{Times New Roman} % 设置英文字体


% =================================
% 数学与公式
% =================================
\usepackage{amsmath}        % AMS 提供的数学公式增强宏包，提供 align, gather 等环境
\usepackage{amssymb}        % 包含大量数学符号，如 R（实数集）、\lesssim 等
\usepackage{amsfonts}       % AMS 提供的数学字体，用于特殊字体（如花体字母）
\usepackage{bm}             % 使公式中的粗体符号（如向量）显示为直立粗体 (\bm)

% =================================
% 图形、表格与浮动体
% =================================
\usepackage{graphicx}       % 引入图片（使用 \includegraphics 命令）
\usepackage{float}          % 增强对浮动体（图、表）的控制，例如 H 定位符
\usepackage{booktabs}       % 绘制专业、美观的三线表（使用 \toprule, \midrule, \bottomrule）
\usepackage{multirow}       % 表格中合并行
\usepackage{makecell}       % 表格中换行和调整单元格边距
\usepackage{caption}        % 更好地自定义图表标题格式
\usepackage{subcaption}     % 在一个浮动体中并排放置多个子图或子表


% =================================
% 链接与排版增强
% =================================
\usepackage{hyperref}       % 生成可点击的超链接（目录、引用、参考文献等），通常应最后加载
\hypersetup{
	colorlinks=true,        % 彩色链接
	linkcolor=black,        % 内部链接颜色
	citecolor=blue,         % 引用链接颜色
	urlcolor=blue           % URL 链接颜色
}
\usepackage{enumerate}      % 增强列表环境（itemize, enumerate），允许自定义标签格式

% =================================
% 算法
% =================================
% 推荐使用下列其中一个宏包：
\usepackage{algorithm}      % 算法浮动体环境
\usepackage{algpseudocode}  % 伪代码排版宏包（与 algorithm 配合使用）
% 或
% \usepackage[linesnumbered, ruled]{algorithm2e} % 另一种流行的算法排版宏包


\usepackage{caption} % 用于定制标题格式
\usepackage{xpatch}  % 用于修改 \captionname 宏（可选，但推荐用于中文/双语环境）

% 1. 定义一个命令来存储中文的“图”和“表”名称，以配合 xeCJK 的中文环境
\renewcommand{\figurename}{图}
\renewcommand{\tablename}{表}

% 2. 使用 caption 宏包定制标题格式：居中对齐，并在编号和文字之间添加空格
\captionsetup{
	justification=centering,   % 居中对齐
	labelsep=quad,             % 标签和文字之间的间距（quad 略大于 space）
	font=normal,               % 字体设置为正常
}


\usepackage{cleveref} % 交叉引用
\crefname{figure}{图}{图}
\Crefname{figure}{图}{图}
\crefname{table}{表}{表}
\Crefname{table}{表}{表}
\crefname{section}{节}{节}
\Crefname{section}{节}{节}
\crefname{equation}{公式}{公式}
\Crefname{equation}{公式}{公式}


\usepackage{cite} % 无自动上标，需要手动上标，下方才用重新定义实现
% =================================
% 重新定义 \cite 命令为上标（带方括号）
% =================================
% 1. 备份原始的 \cite 命令
\let\oldcite\cite 
% 2. 重新定义 \cite 命令，将其内容放入 \textsuperscript{}
\renewcommand*\cite[1]{\textsuperscript{\oldcite{#1}}}

%\usepackage[numbers, super]{natbib} % 无方括号，与cite选用

\usepackage{fancyhdr} % 页眉页脚


% 取消摘要环境自带的标题
\renewcommand{\abstractname}{}
\renewcommand{\absnamepos}{empty}

